\begin{thebibliography}{10}
\small
\setlength{\itemsep}{7pt}

\bibitem{AM}
S.~S.~Abhyankar and T.~T.~Moh,
\emph{Embeddings of the line in the plane},
Journal f\"ur die reine und angewandte Mathematik \textbf{276}
(1975), 148--166.
\href{https://doi.org/10.1515/crll.1975.276.148}{DOI: 10.1515/crll.1975.276.148}.

\bibitem{Suzuki}
M.~Suzuki,
\emph{Propri\'et\'es topologiques des polyn\^omes de deux variables
complexes, et automorphismes alg\'ebriques de l'espace $\C^2$},
Journal of the Mathematical Society of Japan \textbf{26}(2)
(1974), 241--257.
\href{https://doi.org/10.2969/jmsj/02620241}{DOI: 10.2969/jmsj/02620241}.

\bibitem{Abhyankar}
S.~S.~Abhyankar,
\emph{On the semigroup of a meromorphic curve}, Part~I,
in \emph{Proceedings of the International Symposium on Algebraic
Geometry, Kyoto 1977}, Kinokuniya, Tokyo, 1978, 249--414.
Chapter~11, pp.~413--414.
The historical locations are documented in
\href{https://www.numdam.org/item/RSMUP_1985__74__115_0.pdf}{Craighero's
Introduction}, pp.~115--116; the original volume was not consulted here.

\bibitem{Craighero}
P.~C.~Craighero,
\emph{About Abhyankar's conjectures on space lines},
Rendiconti del Seminario Matematico della Universit\`a di Padova
\textbf{74} (1985), 115--122.
\href{https://www.numdam.org/item/RSMUP_1985__74__115_0/}{Journal archive};
\href{https://www.numdam.org/item/RSMUP_1985__74__115_0.pdf}{open PDF}.
See the Introduction and Proposition~3.

\bibitem{Shastri}
A.~R.~Shastri,
\emph{Polynomial representations of knots},
T\^ohoku Mathematical Journal (2) \textbf{44}(1) (1992), 11--17.
\href{https://doi.org/10.2748/tmj/1178227371}{DOI: 10.2748/tmj/1178227371};
\href{https://www.jstage.jst.go.jp/article/tmj1949/44/1/44_1_11/_pdf}{open PDF}.
See \S1 and \S3, Example~1.

\bibitem{KP}
P.-V.~Koseleff and D.~Pecker,
\emph{Chebyshev diagrams for two-bridge knots},
Geometriae Dedicata \textbf{150} (2011), 405--425.
\href{https://doi.org/10.1007/s10711-010-9514-7}{DOI: 10.1007/s10711-010-9514-7};
\href{https://arxiv.org/abs/0909.3281}{arXiv:0909.3281}.
See the Introduction, Definition~1.2 and the identification of $H(3,5,7)$.

\bibitem{KZ}
S.~Kaliman and M.~Zaidenberg,
\emph{Affine modifications and affine hypersurfaces with a very
transitive automorphism group},
Transformation Groups \textbf{4} (1999), 53--95.
\href{https://doi.org/10.1007/BF01236662}{DOI: 10.1007/BF01236662};
\href{https://arxiv.org/abs/math/9801076}{arXiv:math/9801076}.
See Proposition~2.1 and Corollary~2.2.

\bibitem{EV}
E.~Edo and S.~V\'en\'ereau,
\emph{Length 2 variables of $A[x,y]$ and transfer},
Annales Polonici Mathematici \textbf{76} (2001), 67--76.
\href{https://doi.org/10.4064/ap76-1-6}{DOI: 10.4064/ap76-1-6};
\href{https://www.impan.pl/shop/publication/transaction/download/product/85204}{open PDF}.
See Proposition~1, Theorem~4 and \S5.

\bibitem{EVM}
E.~Edo, A.~van den Essen and S.~Maubach,
\emph{A note on $k[z]$-automorphisms in two variables},
Journal of Pure and Applied Algebra \textbf{213} (2009), 1197--1200.
\href{https://doi.org/10.1016/j.jpaa.2008.11.040}{DOI: 10.1016/j.jpaa.2008.11.040};
\href{https://arxiv.org/abs/0809.0767}{arXiv:0809.0767}.
See \S1, Lemma~1.1, which recalls an earlier construction.

\bibitem{BD}
S.~M.~Bhatwadekar and A.~K.~Dutta,
\emph{On residual variables and stably polynomial algebras},
Communications in Algebra \textbf{21}(2) (1993), 635--645.
\href{https://doi.org/10.1080/00927879308824585}{DOI: 10.1080/00927879308824585}.
The residual-variable theorem gives the existence of a companion over
the Noetherian $\Q$-algebra $k[x]$.

\end{thebibliography}
